\section{快速上手}\label{sec:start}

\subsection{记忆模块能为你做什么}

OS Agent 的记忆模块让助手在使用过程中\textbf{越用越懂你}。它由两个互补的记忆库组成。
偏好记忆持续观察你的操作习惯：当你在同类任务中总是选择同一个工具、反复要求某种输出
格式、或明确修改过某项配置时，系统会自动把这些倾向沉淀为偏好，在之后的会话中直接
沿用，不再需要你每次重复交代。知识记忆则是一个可追溯的事实库与模板库：你与助手
确认过的事实（项目参数、联系方式、环境信息）以结构化三元组存储，新旧信息冲突时
自动接替并保留历史版本；一次完整的操作流程（例如「构建—排错—修复」）在积累足够
执行证据并经你确认后，可以固化为可复用的工作流模板。

所有记忆数据都以普通 JSON 文件保存在你的工作区 \code{.amadeus/} 目录下，不依赖任何
数据库或云服务，你可以随时用文本编辑器直接查看、备份或删除。涉及个人敏感信息的
内容（手机号、身份证号、银行卡号、邮箱、密码、API 密钥）在写入记忆前会被自动替换
为占位符；你也可以随时用一句自然语言让助手「忘掉」指定记忆（见第~\ref{sec:forget}~章）。

\subsection{系统要求与安装}

系统要求为 Rust 稳定版工具链（\code{rustup} 安装即可），可选准备一个嵌入服务用于
语义检索（不准备也不影响其余功能，系统会退回内置离线嵌入，见~\S\ref{sec:config-emb}）。
在仓库根目录执行：

\begin{lstlisting}[language=bash]
cargo build --features full
cargo run --features full                # 启动终端交互（TUI）
cargo run --features full -- --server 8080   # 或启动 HTTP 服务
\end{lstlisting}

首次运行后，助手会在当前工作区自动创建 \code{.amadeus/} 目录并初始化全部记忆文件。
各文件的含义见表~\ref{tab:files}，日常使用中你通常不需要直接修改它们。

\begin{table}[htbp]
  \centering\small
  \caption{记忆数据文件一览（\code{.amadeus/} 目录）}
  \label{tab:files}
  \begin{tabular}{l l}
    \toprule
    文件 & 内容 \\
    \midrule
    \code{preferences.json} & 偏好记忆：学到的偏好及其历史版本 \\
    \code{structured_memory.json} & 知识记忆：事实、关联、模板 \\
    \code{mid_term_memory.json} & 中期记忆：跨会话的任务与决策记录 \\
    \code{memory.json} & 长期事实记忆（注入对话上下文） \\
    \code{rag_index.json} & 语义检索的向量索引 \\
    \bottomrule
  \end{tabular}
\end{table}

在银河麒麟桌面操作系统上的部署要点与适配说明见~\S\ref{sec:kylin-deploy}。

\subsection{五分钟体验记忆闭环}

安装启动后，你可以按下面的顺序快速体验记忆模块的核心能力。

第一步，\textbf{让系统学到偏好}。用自然语言让助手完成几个同类任务并指定同一个工具，
例如连续两次要求「用 ripgrep 统计 xxx」。系统在工具执行成功且样本充足后自动生成
工具选择偏好，之后同类任务会优先使用该工具。你也可以直接告诉助手「以后输出都用
简洁格式」，明确的指令会立即成为高置信度偏好。

第二步，\textbf{沉淀一条知识}。对助手说「记住：我的项目构建参数是 xxx」。助手会调用
\code{kylin\_memory} 工具把这条信息写成结构化事实。之后你随时可以问「我的构建参数
是什么」，助手会从知识库检索作答。

第三步，\textbf{验证冲突接替}。接着说「更新一下：构建参数改为 yyy」。系统不会新旧
混杂，而是把旧值标记为历史版本、以新值接替；你可以要求助手「查一下这条事实的修改
历史」看到完整的版本链。

第四步，\textbf{体验精准遗忘}。说「忘掉我的构建参数，但保留项目名称」。助手会先给出
一份删除预览（列出将删除的内容与保留的内容），等你明确确认后才真正删除，并在删除后
自动校验磁盘上已无残留。
